% =====================================================================
\chapter{Радиомодули: ExpressLRS и мультипротокол}
\label{ch:rf}
% =====================================================================
\chapterintro
  {как включить внутренний или внешний модуль, привязать приёмник ExpressLRS или FrSky,
   что означают частота пакетов, соотношение телеметрии, режимы каналов, мощность и failsafe}
  {приёмник с совместимым протоколом, модель без пропеллеров}

\section{Какие бывают модули}

\begin{center}\small
\renewcommand{\arraystretch}{1.2}
\begin{tabularx}{\linewidth}{@{}p{30mm}p{30mm}X@{}}
\toprule
Модуль & Режим в EdgeTX & Совместимые приёмники \\
\midrule
ExpressLRS 2,4\,ГГц & \menu{CRSF} & приёмники ELRS 2,4\,ГГц той же мажорной версии \\
4-в-1 мультипротокол & \menu{MULTI} & FrSky D8/D16/ACCST, Flysky AFHDS/AFHDS2A, Spektrum DSM2/DSMX,
  Futaba S-FHSS, Hubsan, Syma и десятки других (чипы CC2500, A7105, CYRF6936, NRF24L01) \\
CC2500 & \menu{MULTI} & только протоколы на чипе CC2500: FrSky D8/D16, Futaba S-FHSS,
  Corona, Hitec, Redpine и др. \\
Внешний TBS Crossfire / ELRS & \menu{CRSF} (External RF) & Crossfire / ELRS \\
Внешний PPM-модуль & \menu{PPM} & любой модуль с PPM-входом \\
\bottomrule
\end{tabularx}
\end{center}

\section{Включение модуля в модели}

В \mpath{MDL\sep Setup} найдите раздел \menu{Internal RF} (или \menu{External RF} для
модуля в отсеке JR) и выберите режим:

\begin{itemize}
  \item \menu{Mode} = \menu{CRSF} для ELRS или \menu{MULTI} для мультипротокола;
  \item \menu{Receiver} (\menu{Rx number}) --- номер приёмника для функции Model Match;
  \item \menu{Channel range} --- какие каналы передавать (например, \sw{CH1}--\sw{CH16});
  \item \menu{Baudrate} (для CRSF) --- скорость связи пульта с модулем.
\end{itemize}

\begin{note}
Одновременно можно включить и внутренний, и внешний модуль (например, для двух разных
приёмников), но обычно используют один. Если вы вставили внешний модуль, переключите
\menu{Internal RF} в \menu{OFF} --- это снизит энергопотребление и помехи.
\end{note}

\section{ExpressLRS}

ExpressLRS (ELRS) --- открытая радиосистема с большой дальностью и малой задержкой, стандарт
де-факто в FPV. EdgeTX общается с модулем по протоколу CRSF, а настраивается модуль
Lua-скриптом.

\subsection{Скрипт ExpressLRS}

\mpath{SYS\sep Tools\sep ExpressLRS} (файл \code{SCRIPTS/TOOLS/elrsV3.lua}). Если скрипта нет
в списке, обновите содержимое SD-карты или скачайте его из ExpressLRS Configurator.

\begin{center}
\lcd{\inv{ExpressLRS}\hfill 0/250\\
Packet Rate\hfill 250Hz(-108dbm)\\
Telem Ratio\hfill Std (1:64)\\
Switch Mode\hfill Wide\\
Model Match\hfill Off\\
TX Power (250mW)\hfill >\\
VTX Administrator\hfill >\\
WiFi Connectivity\hfill >\\
{[Bind]}}
\end{center}

Число \code{0/250} в заголовке --- счётчик плохих/хороших пакетов; если там \code{--},
скрипт не видит модуль (проверьте \menu{Internal RF} = \menu{CRSF}).

\subsection{Привязка (binding)}

\textbf{Способ 1 --- фраза привязки.} Если и в передатчик, и в приёмник при прошивке
записана одна и та же binding phrase, они связываются \emph{автоматически}, никаких действий
не нужно. Это самый удобный способ для группы: у каждого ученика --- своя фраза.

\textbf{Способ 2 --- режим привязки.}
\begin{enumerate}
  \item Переведите приёмник в режим привязки: трижды включите и выключите его питание
        (на третий раз оставьте включённым) --- светодиод начнёт быстро мигать двойными
        вспышками. У приёмников ELRS 3.x можно также зайти в web-интерфейс приёмника или
        нажать кнопку bind, если она есть.
  \item В скрипте ExpressLRS выберите \menu{[Bind]}.
  \item Светодиод приёмника должен загореться постоянно, в скрипте появится счётчик
        хороших пакетов, а на пульте --- голосовое «Telemetry recovered».
\end{enumerate}

\begin{caution}
Если в передатчик зашита фраза привязки, а приёмник привязан через режим bind, то после
смены фразы на передатчике связь пропадёт. Придерживайтесь одного способа.
\end{caution}

\subsection{Частота пакетов (Packet Rate)}

\begin{center}\small
\renewcommand{\arraystretch}{1.15}
\begin{tabularx}{\linewidth}{@{}p{30mm}p{28mm}X@{}}
\toprule
Режим & Чувствительность & Когда использовать \\
\midrule
50\,Гц & около $-115$\,дБм & максимальная дальность, самолёты/крылья «на дальняк» \\
100/150\,Гц & около $-112$\,дБм & дальние полёты, крылья \\
250\,Гц & около $-108$\,дБм & универсальный режим для FPV-фристайла \\
333\,Гц Full, 500\,Гц & около $-105$\,дБм & гонки, быстрый отклик \\
F500, F1000, K1000 & хуже & гонки на короткой трассе \\
D250, D500 & & режимы с повтором пакетов (DVDA) \\
\bottomrule
\end{tabularx}
\end{center}

Чем выше частота --- тем меньше задержка, но тем меньше чувствительность и дальность.
Для учебных полётов и фристайла 250\,Гц --- хороший выбор. Режимы 1000\,Гц требуют высокой
скорости связи пульта с модулем (\menu{Baudrate} от 921\,600 или 1,87\,Мбод) --- если
режим недоступен в списке, увеличьте \menu{Baudrate} в настройках модуля.

\subsection{Соотношение телеметрии (Telem Ratio)}

Доля пакетов, отводимых под телеметрию: \menu{Std} (стандарт для выбранной частоты),
\menu{Off}, \menu{1:128} … \menu{1:2}, \menu{Race} (телеметрия отключается в арме).
Больше телеметрии --- выше частота обновления данных (например, для Lua-скриптов
Betaflight), но меньше пакетов управления.

\subsection{Режимы каналов (Switch Mode)}

ELRS передаёт каналы \sw{CH1}--\sw{CH4} (стики) в полном разрешении каждый пакет, а
«вспомогательные» каналы --- по-разному:

\begin{description}[font=\sffamily\bfseries\color{etx}]
  \item[Hybrid] \sw{CH5} (AUX1) --- 1 бит (два положения) в каждом пакете; \sw{CH6}--\sw{CH11}
        по очереди, с низким разрешением (6 положений); \sw{CH12} --- 16 положений.
  \item[Wide] \sw{CH5} --- 1 бит; \sw{CH6}--\sw{CH12} по очереди, 7 бит (128 положений) ---
        удобно, если на AUX-каналах крутилки (подвес камеры, свет).
  \item[Full Res] (в режимах вроде 333\,Full и K1000\,Full) --- 8, 12 или 16 каналов
        в полном разрешении.
\end{description}

\begin{caution}
В ExpressLRS канал \sw{CH5} (AUX1) используется как \textbf{сигнал арма}: он передаётся
всегда и в обоих положениях влияет на поведение системы (например, мощность и телеметрию).
Назначайте на \sw{CH5} \emph{двухпозиционный} переключатель арма и не используйте \sw{CH5}
для других функций.
\end{caution}

\subsection{Мощность (TX Power)}

\begin{itemize}
  \item \menu{Max Power} --- максимальная мощность: 10, 25, 50, 100, 250\,мВт (TX12);
        до 500\,мВт и 1000\,мВт (Boxer).
  \item \menu{Dynamic} --- динамическая мощность: модуль сам снижает мощность, пока связь
        хорошая, и поднимает её при ухудшении. Может управляться переключателем (\menu{AUX9}
        и т.\,п.).
  \item \menu{Fan Thresh} (Boxer) --- с какой мощности включать вентилятор.
\end{itemize}

\begin{tip}
Мощность 1\,Вт на Boxer заметно разряжает аккумулятор и нагревает модуль. Для полётов в
пределах видимости достаточно 100--250\,мВт с включённой динамической мощностью.
Соблюдайте ограничения мощности, действующие в вашей стране.
\end{tip}

\subsection{Model Match}

Если \menu{Model Match} включён, приёмник запоминает номер модели (\menu{Receiver} на
странице \menu{Setup}) и не будет работать, если в пульте выбрана другая модель. Это защищает
от ситуации «выбрал модель самолёта, а подключился к квадрокоптеру». Порядок: задайте
уникальный \menu{Receiver} в каждой модели → привяжите приёмник → включите Model Match.

\subsection{Failsafe в ELRS}

При потере связи приёмник ELRS по умолчанию \emph{перестаёт выдавать сигнал} --- полётный
контроллер это распознаёт и включает свой failsafe (падение или возврат домой по GPS).
Для PWM-приёмников (с сервами напрямую) положения failsafe задаются в web-интерфейсе
приёмника или командой \menu{Set failsafe} (зависит от версии ELRS).

\section{Мультипротокольный модуль (4-в-1, CC2500)}

\subsection{Настройка протокола}

В \mpath{MDL\sep Setup\sep Internal RF} выберите \menu{Mode} = \menu{MULTI}, затем:

\begin{description}[font=\sffamily\bfseries\color{etx}]
  \item[Protocol] например \menu{FrSky X} (D16), \menu{FrSky D} (D8), \menu{FrSkyX2}
        (ACCST~D16~v2), \menu{FlySky 2A}, \menu{DSM}.
  \item[Subtype] подтип: \menu{D16}, \menu{D16 8ch}, \menu{LBT(EU)}, \menu{EU LBT 8ch} и т.\,д.
        Для FrSky важно совпадение версии прошивки приёмника: ACCST~v1 или v2, FCC или EU-LBT.
  \item[RF Freq. fine tune] подстройка частоты. Если связь рвётся на расстоянии пары метров,
        попробуйте значения от $-40$ до $+40$.
  \item[Receiver No.] номер приёмника (аналог Model Match).
  \item[Failsafe] поведение при потере связи (см. ниже).
  \item[Bind / Range] привязка и проверка дальности.
\end{description}

\subsection{Привязка FrSky D16}

\begin{enumerate}
  \item Выберите протокол и подтип так, чтобы они совпадали с прошивкой приёмника.
  \item Включите приёмник, удерживая на нём кнопку \code{F/S} (bind) --- светодиоды начнут мигать.
  \item На пульте выберите \menu{Bind}. Пульт начнёт пищать.
  \item Через пару секунд светодиод приёмника изменит режим мигания. Остановите привязку
        на пульте и перезапустите питание приёмника.
  \item Приёмник должен светиться постоянно (обычно зелёным) --- связь есть.
\end{enumerate}

\subsection{Failsafe}

\begin{center}\small
\begin{tabularx}{\linewidth}{@{}p{30mm}X@{}}
\toprule
\menu{Not set} / \menu{Receiver} & используется то, что записано в приёмнике \\
\menu{Hold} & удерживать последние значения (опасно для мотора!) \\
\menu{Custom} & задать положения каналов вручную (газ --- минимум) \\
\menu{No pulses} & прекратить выдачу сигнала; нужен полётному контроллеру для распознавания
  потери связи \\
\bottomrule
\end{tabularx}
\end{center}

\begin{caution}
Всегда проверяйте failsafe на земле со снятыми пропеллерами: включите модель, выключите
пульт и посмотрите, что делают мотор и сервы (или что показывает полётный контроллер).
Мотор должен остановиться.
\end{caution}

\section{Проверка дальности}

\menu{Range} (для Multi) уменьшает мощность передатчика; при этом связь должна сохраняться
на расстоянии около 30--40\,м. Для ELRS проверка делается по телеметрии: отойдите от модели и
посмотрите значения \sw{1RSS} и \sw{RQly} (глава~\ref{ch:telemetry}).

\begin{summary}
\begin{itemize}
  \item ELRS = режим \menu{CRSF} + Lua-скрипт ExpressLRS; Multi = режим \menu{MULTI} +
        протокол/подтип.
  \item ELRS связывается по фразе привязки или через \menu{[Bind]}.
  \item \sw{CH5} в ELRS --- канал арма, только двухпозиционный переключатель.
  \item Failsafe проверяется до первого полёта. Мотор при потере связи должен останавливаться.
\end{itemize}
\end{summary}
